% chapters/ch01.tex —— 第 1 章
% 体例：具体失败 → 设计问题 → 机制解释 → 对照实验 → 条件变化 → 选择原则
% 数字身份纪律：外部数值使用 \cited{}；协议配额与本书状态另明示身份。
\chapter{从自主执行到自进化}
\label{ch01}

\chaptergoal{一次任务结束后，系统留下了什么？}
\chapterprops{P1、P1b（第3章的预算与只读对照）；本章为后续研究问题提供术语与比较依据}
\chapterfig{图~\ref{fig:timescale}：三层时间尺度上的状态保留范围}
\chapterpitfall{按 rollout 配额更新 harness 后，v2 的两模型留出均值仅增 0.6 个百分点}

% ---------- 1. 具体失败 ----------
\begin{failurebox}
2026 年一项按统一 rollout 协议比较的方法研究（本文此处采用2026年8月27日的v2）给出了一个值得正视的结果\cite{harness-eval-2607}。
研究者在 Terminal-Bench 2.1 上比较四条路径：并行采样、顺序精化、harness evolution、harness scaling。
v2的§4.1设定计算配额$K=\cited{5}$，harness evolution每任务每轮采样$m=\cited{1}$条轨迹，单次生成上限为\cited{128k} token；这些是不同单位，不能合称“修改轮数”。四条路径按论文规定的 rollout 数量与生成长度上限比较。这个协议约束了采样机会，
却不能直接推出实际 token、墙钟或货币成本逐项相等；读表时要保留这一预算单位。

在没有单元测试反馈时，初始 harness 直接采样的 pass@1 平均是 \cited{68.2}。
顺序精化把它提到 \cited{69.3}，harness scaling 提到 \cited{71.8}，而 harness evolution 落在 \cited{67.4}——\textbf{低于不做任何演化的起点}。
同一批实验里，并行采样的三模型平均点估计最高，为\cited{72.3}。

更值得记的是迁移那一栏。研究者把 Terminal-Bench 2.1 切成 45 个训练任务、10 个验证任务、34 个 held-out 测试任务，
在训练集上跑 harness evolution，按验证集挑最好的 harness，再到测试集上测。
结果是：Claude Opus 4.6 涨 \cited{1.2} 个百分点，GPT-5.4 涨 \cited{0.0} 个百分点，平均 \cited{0.6}。
原作者的解释是——进化过程中发现的修改编码的是\textbf{针对具体任务的捷径}，不是更好的 harness 设计原则。
\end{failurebox}

\designquestion{一次系统改动让分数上去了，怎么知道它带来了能力，还是只是一次更贵的重试？}

第3章将展开这个问题的比较方法；在第1章，先把它的对象看清：
上面那组数字里，harness 确实被改了，持久物确实存在，也确实被带到了下一个任务——
在该论文声明的rollout配额下，把机会用于“多采样几条轨迹再挑一条”，平均表现更高。
于是“系统有没有留下东西”和“留下东西有没有用”必须是两个问题，本书从头到尾把它们分开。

% ---------- 2. 设计问题 → 3. 机制解释 ----------
\section{自主、适应、学习、演化}

四个词经常被当成一回事，它们的差别决定了后面每一章在讲什么。

\begin{itemize}
  \item \textbf{自主执行}：系统能独立走完一个动作序列。执行机制可以始终不变。
  \item \textbf{适应}：系统依据当前条件调整行为，例如利用上下文中的示例改变解题方式。调整是否有效，需要另行比较。
  \item \textbf{学习}：经验引起行为或表示的改变。文献中的学习也包含上下文学习，不必都伴随参数更新或跨任务保存。
  \item \textbf{演化}：改变落在\textbf{持久组件}上，并且改变之后下一个任务从变化过的系统出发。
\end{itemize}

前三者都可以在一次任务内部发生；本书对第四件事增加跨任务留存的要求。
这是一组用于组织本书实验的操作性定义，并非给整个学习领域重新划界。社区清单的术语整理提供了相关参照（表~\ref{tab:rsi-terms}）\cite{awesome-rsi}。

用一个具体场景把四个词分开。一个编程 Agent 被要求修复一个仓库里的失败测试：

\begin{itemize}
  \item 它自己读代码、改文件、跑测试、再改，全程没有人类插手——这是\textbf{自主执行}。
  \item 在同一任务中，它修复的第 8 个 bug 比第 1 个快，因为上下文里已经包含前 7 次的报错——这是\textbf{适应}，
        若下一任务清空上下文，这份由临时信息带来的优势便未必保留。
  \item 它把"本项目的测试入口是 \texttt{pytest -q tests/}"写进了一个跨任务保留的文件——这是一个经验更新候选，仍需接通读取路径。
  \item 它把这条经验写进持久记忆，并且\textbf{在下一批任务里被检索到、改变了行为}，并能接受独立检验——这构成了本书要研究的\textbf{演化过程}；是否有益另行判定。
\end{itemize}

第三件和第四件之间隔着一整套工程：写入、组织、检索、验证、发布、回滚。本书讲的正是这一整套。

\begin{table}[htbp]
\centering
\begin{tabular}{@{}p{3.6cm}p{5.0cm}p{2.4cm}p{2.6cm}@{}}
\toprule
英文术语 & 定义 & 本书时间尺度 & 对应组件 \\
\midrule
self-refinement & 改善当前输出，系统无持久变化 & episode $e$ 及以下 & — \\
persistent self-improvement & 权重/记忆/技能/提示/运行脚手架改变并带入下一轮 & version $k$ & $\text{Updatable}\subseteq\{\theta,P,M,L,G\}$ \\
recursive self-improvement (RSI) & 产生改进的机制本身也成为被改进对象 & version $k$ 上的二层结构 & 第 12 章的修改器 \\
self-evolving（本书术语） & 修改持久组件，且该变化须可独立检验 & version $k$，且变化须对 $\mathcal V$ 可测 & $\text{Updatable}\subseteq\{\theta,P,M,L,G\}$，$\mathcal V$ 冻结 \\
\bottomrule
\end{tabular}
\caption{四词并排。前三项参照社区清单 \texttt{lobehub/awesome-rsi} 的术语说明，第四项是本书用法；这不是学界统一标准，也不是互斥类别。}
\label{tab:rsi-terms}
\end{table}

四行回答的不是同一个问题，需要分别检查：
“我们的上下文变长，所以发生了跨任务改进”——长度本身不能说明保留边界；
``我们有持久物，所以它是自进化的''——持久物存在不等于持久物被修改，只读不写停在第二行之外；
``我们改了权重，所以是递归自我改进''——第三行要求被改的是产生改进的那个机制，不是机制产出的东西；
``它一直在变强''——第四行要的是可独立检验的变化，能力无上界那一部分本书整体不声称（第 16 章）。

\begin{definition}[自进化]
Agent 根据经验与反馈修改\textbf{持久组件}，使后续任务从一个已经发生变化的系统出发；这种变化是否有效，需要独立检验。
\end{definition}

定义与效果判决要分开。经验已经改变后续任务使用的组件，却还没有效果证据，可以记录为效果未明的更新；不能因为成绩未知就否认更新发生过。反过来，只有测量而没有持久更新，也不能据此声称发生了本书意义上的自进化。

因此，“上下文很长”不足以支持跨任务改进的主张。
若上下文仅在任务内部保留，结束即清空，且没有其他持久写入，下一个任务仍从同一个 $A_k$ 出发。
这个判断只涉及本书的保留边界，不否认模型在当前上下文中利用示例与反馈进行学习。GPT-3论文的少样本设置就明确使用文本示例而不在这些任务上做梯度更新或微调\cite{brown-few-shot-2005}。
要让这段经验成为本书研究的跨任务更新，需要把相应改变保存到 $M$、$P$、$L$、$G$ 或 $\theta$，并接通后续使用路径。
保存之后还应回答如何选择、加载与撤回；具体需要检索、淘汰还是重训，取决于更新对象，不能给五种组件规定同一套操作。

\section{三层时间尺度}
\label{sec:timescale}

一个 Agent 系统里发生变化的东西，按\textbf{保留多久}分成三层（表~\ref{tab:scale}）。

\begin{table}[htbp]
\centering
\begin{tabular}{@{}p{1.6cm}p{1.4cm}p{5.0cm}p{4.6cm}@{}}
\toprule
层 & 变量 & 变化的东西 & 典型混淆 \\
\midrule
step $t$ & $H_{k,e,t}$ & 工作上下文与当前输入条件 & 长度增长不证明持久更新 \\
episode $e$ & $\tau_{k,e}$ & 单次任务内反思与重试 & 自精炼只在这一层 \\
version $k$ & $A_k$ & 持久组件 & 跨 episode 写记忆 \\
\bottomrule
\end{tabular}
\caption{三层时间尺度。只有 version 层的变化算自进化。}
\label{tab:scale}
\end{table}

% figures/timescale.tex —— 第 1 章核心图（全书引用入口）
% 用法：在 book/chapters/ch01.tex 中 % figures/timescale.tex —— 第 1 章核心图（全书引用入口）
% 用法：在 book/chapters/ch01.tex 中 % figures/timescale.tex —— 第 1 章核心图（全书引用入口）
% 用法：在 book/chapters/ch01.tex 中 \input{figures/timescale}
% 图主体在 timescale-body.tex，便于单独编译与复用。
\begin{figure}[htbp]
\centering
\input{figures/timescale-body}
\caption{三层时间尺度上的状态保留范围。临时执行状态按任务协议复位，轨迹可独立归档；虚线表示按更新规则提炼并接入持久组件的可能路径，不表示每份归档都会被采用。只有进入version层并被后续任务继承、使$A_{k+1}\neq A_k$的变化，才属于本书的自进化研究对象。}
\label{fig:timescale}
\end{figure}

% 图主体在 timescale-body.tex，便于单独编译与复用。
\begin{figure}[htbp]
\centering
% figures/timescale-body.tex —— 三层时间尺度图的纯 tikz 主体
% 单独编译入口：figures/timescale-standalone.tex
% 全书引用入口：figures/timescale.tex（含 figure 环境与编号）
% 图内不含任何数字，因此不受 \cited{} / \calcnum{} / \measured{} 身份纪律约束。
% version 行画出完整五类持久组件（$\theta,P,M,L,G$），与式 (1.1) 的 Updatable 集合一致。
\begin{tikzpicture}[
  box/.style={draw, rounded corners=3pt, minimum height=9mm, minimum width=18mm, font=\small},
  lbl/.style={font=\small, anchor=east},
  note/.style={font=\footnotesize, anchor=west, gray!60!black}
]
  % ---- step t ----
  \node[lbl] at (1.4,2.4) {step $t$};
  \node[box, fill=gray!10] (s1) at (3.0,2.4) {消息};
  \node[box, fill=gray!10] (s2) at (6.0,2.4) {工具结果};
  \node[box, fill=gray!10] (s3) at (9.0,2.4) {消息};

  % ---- episode e ----
  \node[lbl] at (1.4,1.2) {episode $e$};
  \node[box, fill=blue!8] (e1) at (4.0,1.2) {临时工作状态};
  \node[box, fill=blue!8] (e2) at (9.0,1.2) {轨迹可另存};

  % ---- version k：五类持久组件 ----
  \node[lbl] at (1.4,0.0) {version $k$};
  \node[box, fill=orange!12] (k1) at (3.0,0.0) {$M$ 记忆};
  \node[box, fill=orange!12] (k2) at (5.0,0.0) {$P$ 指令};
  \node[box, fill=orange!12] (k3) at (7.0,0.0) {$L$ 技能};
  \node[box, fill=orange!12] (k4) at (9.0,0.0) {$G$ 工作流};
  \node[box, fill=orange!12] (k5) at (11.0,0.0) {$\theta$ 参数};

  % ---- 层内与层间关系 ----
  \draw[-{Stealth[length=2mm]}] (s1)--(s2);
  \draw[-{Stealth[length=2mm]}] (s2)--(s3);
  \draw[-{Stealth[length=2mm]}, dashed] (e2.south) -- (k3.north);

  % ---- 保留范围的标注 ----
  \draw[gray!60!black, line width=0.3pt] (1.8,1.75) -- (12.4,1.75);
  % Reset and archive conditions are stated in the caption; keeping them here
  % would collide with the step-row boxes at the final printed size.
  \draw[gray!60!black, line width=0.3pt] (1.8,-0.55) -- (12.4,-0.55);
  \node[note] at (1.8,-0.95) {跨任务保留：只有这里的变化满足 $A_{k+1}\neq A_k$};
\end{tikzpicture}

\caption{三层时间尺度上的状态保留范围。临时执行状态按任务协议复位，轨迹可独立归档；虚线表示按更新规则提炼并接入持久组件的可能路径，不表示每份归档都会被采用。只有进入version层并被后续任务继承、使$A_{k+1}\neq A_k$的变化，才属于本书的自进化研究对象。}
\label{fig:timescale}
\end{figure}

% 图主体在 timescale-body.tex，便于单独编译与复用。
\begin{figure}[htbp]
\centering
% figures/timescale-body.tex —— 三层时间尺度图的纯 tikz 主体
% 单独编译入口：figures/timescale-standalone.tex
% 全书引用入口：figures/timescale.tex（含 figure 环境与编号）
% 图内不含任何数字，因此不受 \cited{} / \calcnum{} / \measured{} 身份纪律约束。
% version 行画出完整五类持久组件（$\theta,P,M,L,G$），与式 (1.1) 的 Updatable 集合一致。
\begin{tikzpicture}[
  box/.style={draw, rounded corners=3pt, minimum height=9mm, minimum width=18mm, font=\small},
  lbl/.style={font=\small, anchor=east},
  note/.style={font=\footnotesize, anchor=west, gray!60!black}
]
  % ---- step t ----
  \node[lbl] at (1.4,2.4) {step $t$};
  \node[box, fill=gray!10] (s1) at (3.0,2.4) {消息};
  \node[box, fill=gray!10] (s2) at (6.0,2.4) {工具结果};
  \node[box, fill=gray!10] (s3) at (9.0,2.4) {消息};

  % ---- episode e ----
  \node[lbl] at (1.4,1.2) {episode $e$};
  \node[box, fill=blue!8] (e1) at (4.0,1.2) {临时工作状态};
  \node[box, fill=blue!8] (e2) at (9.0,1.2) {轨迹可另存};

  % ---- version k：五类持久组件 ----
  \node[lbl] at (1.4,0.0) {version $k$};
  \node[box, fill=orange!12] (k1) at (3.0,0.0) {$M$ 记忆};
  \node[box, fill=orange!12] (k2) at (5.0,0.0) {$P$ 指令};
  \node[box, fill=orange!12] (k3) at (7.0,0.0) {$L$ 技能};
  \node[box, fill=orange!12] (k4) at (9.0,0.0) {$G$ 工作流};
  \node[box, fill=orange!12] (k5) at (11.0,0.0) {$\theta$ 参数};

  % ---- 层内与层间关系 ----
  \draw[-{Stealth[length=2mm]}] (s1)--(s2);
  \draw[-{Stealth[length=2mm]}] (s2)--(s3);
  \draw[-{Stealth[length=2mm]}, dashed] (e2.south) -- (k3.north);

  % ---- 保留范围的标注 ----
  \draw[gray!60!black, line width=0.3pt] (1.8,1.75) -- (12.4,1.75);
  % Reset and archive conditions are stated in the caption; keeping them here
  % would collide with the step-row boxes at the final printed size.
  \draw[gray!60!black, line width=0.3pt] (1.8,-0.55) -- (12.4,-0.55);
  \node[note] at (1.8,-0.95) {跨任务保留：只有这里的变化满足 $A_{k+1}\neq A_k$};
\end{tikzpicture}

\caption{三层时间尺度上的状态保留范围。临时执行状态按任务协议复位，轨迹可独立归档；虚线表示按更新规则提炼并接入持久组件的可能路径，不表示每份归档都会被采用。只有进入version层并被后续任务继承、使$A_{k+1}\neq A_k$的变化，才属于本书的自进化研究对象。}
\label{fig:timescale}
\end{figure}


\begin{definition}[自进化判据]
经验与反馈驱动的更新，需要改变后续任务实际加载的持久组件，即在明确组件语义下有 $A_{k+1}\neq A_k$，并能够独立检验。单纯文件差异不是充分条件，更新有益也不是定义预先保证的结果。
\end{definition}

持久性要求有文献参照：一篇综述第2.1节把持久的策略改变列为纳入条件，同时说明其讨论范围包含较早期的原型演化\cite{self-evolving-survey-2507}；另一篇第3.1节把临时执行状态 $X_t$ 与模型及运行结构分开\cite{self-improvements-survey-2607}。
这些来源支持区分变化对象与保留范围，但并未逐字规定本书的任务边界。本书进一步把后续任务实际加载的组件作为检查对象，不把这一操作化说成所有综述共用的排他定义。

\begin{sloppypar}
不能只按方法名称分层。ReAct~\cite{react-2210} 的推理与行动轨迹、Self-Refine~\cite{self-refine-2303} 的输出修订，在任务结束后清空的实现中属于临时状态。
Reflexion~\cite{reflexion-2303} 则明确保留反思供后续尝试读取。若这些尝试都在求解同一道题，它能支持重试中的经验利用，却不能直接证明向新任务迁移。
如果另一种实现把反思带入新任务，就应检验那条持久链路，不能因沿用某个方法名称而一概排除。
本书按状态实际跨过的边界判断，不把论文中的 trial、episode 与本书任务边界无条件画等号。
\end{sloppypar}

判据要能执行才算判据。在本书的实验台里，三层对应三个可以被检查的落点：
上下文只活在一次运行的进程里；轨迹写进 \texttt{trace.json}，留下来了，但下一个任务不会去读它；
只有进入持久组件、并被后续任务实际加载的变化，才构成本书检查的跨任务更新。
于是一个方法落在哪一层，需要核对实际写入与后续加载路径，不能只按方法名称判断。
格 1 的静态基线在每一次 \texttt{run\_episode} 之后检查轨迹中的持久写入字段，非空就抛出异常。
这只能约束已登记的写入路径；字段为空不能证明进程没有通过其他渠道改动文件或外部状态，第2章与附录B进一步说明这一限制。

落点还有一个反直觉的区别：同一段文本，留在 episode 层是任务内反思，
写进 $M$ 就成为持久物；它是否改变后续任务的起点，还要核对加载与使用。内容一模一样，保留范围与作用证据仍可能不同。
这也是为什么“我们记录了很多东西”不能作为自进化的证据：记录落在 $\tau_{k,e}$，判决只看 $A_k$。

\section{系统里会变的那五样东西}
\label{sec:components}

把系统版本写成

\begin{equation}
A_k=(\theta_k,\ P_k,\ M_k,\ L_k,\ G_k\ ;\ \mathcal V_k),
\qquad \text{Updatable}\subseteq\{\theta,P,M,L,G\}
\label{eq:ak}
\end{equation}

五类持久组件的分工见表~\ref{tab:components}。

\begin{table}[htbp]
\centering
\begin{tabular}{@{}p{1.2cm}p{4.4cm}p{4.0cm}p{3.4cm}@{}}
\toprule
符号 & 是什么 & 典型更新方式 & 更新代价的性状 \\
\midrule
$\theta$ & 模型参数 & 监督微调、强化学习 & 训练成本可能较高，逐条经验难定位 \\
$P$ & 指令与持久策略 & 提示搜索、指令演化 & 需计搜索与复测；文件可回退 \\
$M$ & 经验记忆 & 写入、检索、淘汰、重写 & 需计索引、检索与噪声筛查 \\
$L$ & 工具与技能库 & 程序生成、验证后入库 & 需检查执行行为与依赖 \\
$G$ & 工作流与角色结构 & 结构搜索、拓扑改写 & 需计搜索与受影响路径复测 \\
\bottomrule
\end{tabular}
\caption{五类持久组件。$\mathcal V$ 不在其中，原因见本节末尾。}
\label{tab:components}
\end{table}

$\mathcal V$ 是验证器：任务划分、验收器与最终测试集的合称。它被单列在分号右边，因为它\textbf{默认冻结}。

\begin{definition}[三条公理]
\begin{enumerate}
  \item $\mathcal V$ 默认冻结，不在 Updatable 内；当它被纳入 Updatable 时必须单列实验组并说明后果。
  \item Updatable 的每一次扩大，都要重新回答“凭什么认为有用”。
  \item 只有进入 version 层（$k$）的变化才叫自进化。
\end{enumerate}
\end{definition}

公理 1 不是洁癖。一个系统如果既能改自己的实现，又能改测试，还能改评分逻辑，
那么它的分数上涨就不足以单独说明原目标上的能力提高，还需独立参照与新旧评价的对应。第14章会把这一类失效拆开讲；第12章的元进化案例（\cite{dgm-2505}）
会给出一个更细的样本：隐藏的检查函数未必被改动，提供给它的运行日志却可能被改变。

公理1在本书实验台上只得到部分实现。任务集v1的验收测试被列入指定写工具的保护集合，通常的直接写请求会被拒绝；但保护集合不是操作系统权限。旧路径检查还存在相邻同名前缀目录误放行，普通子进程也能绕过写工具，附录B给出只使用临时虚构文件的复现。题目成绩由\texttt{task.verify}执行验收命令产生，不由模型自评；这确定了判定来源，却不能单独证明验收程序未被其他执行路径影响。完整隔离仍是新确认性运行需要补齐的条件。

公理 2 是本书采用的更新范围约定，\texttt{spec/glossary.md} 第 3 节登记了它的本书用法。
这项登记不构成“文献中没有对应做法”的新颖性证明。既然用它组织后续检验，就要附一份能照着做的说明：
Updatable 每扩大一档，系统多得到一项能力，同时多背上一类代价。
原有检验是否覆盖新增代价，需要逐项核对（表~\ref{tab:updatable}）。

\begin{table}[htbp]
\centering
\small
% 窄列内的中日韩文字可拉伸空间极小，两端对齐会逐行产生 Underfull；本表各列改为左对齐。
\begin{tabular}{@{}>{\raggedright\arraybackslash}p{1.4cm}>{\raggedright\arraybackslash}p{2.8cm}>{\raggedright\arraybackslash}p{4.4cm}>{\raggedright\arraybackslash}p{4.6cm}@{}}
\toprule
新纳入的组件 & 多得到什么 & 多付出什么 & 与主张对应的检验 \\
\midrule
$P$ & 能改自己的指令与持久策略 & 编辑文本可以很快，但搜索、复测与外部效果补偿未必便宜 & 在声明的读取者上比较提示效果。若更换读取者，应重新评价。同一提示的解释与效果可能随模型变化，恢复旧文本也不撤销已经发生的外部动作 \\
$M$ & 能跨任务积累经验条目 & 写入、索引与检索均有成本；保留规则决定哪些经验继续可见 & 要识别继续更新的作用，应设置来源一致的只读持久对照（P1b、格2），区分已有经验与后续更新的贡献 \\
$L$ & 能把做法固化成可复用的能力 & 需生成、检查并维护技能与依赖 & 明确准入信号及其覆盖范围；模型判断成功与可执行断言通过要分开记录，复用后再检验适用范围和成本（P4） \\
$G$ & 能改任务被分解与分派的方式 & 结构搜索和受影响路径的检查可能扩大成本 & 同时报告任务回归与资源口径。AFlow已做跨模型比较，但测试执行成本曲线不能单独回答全流程回本问题 \\
$\theta$ & 能改系统的默认行为本身 & 训练成本可能较高，逐条经验难定位 & 需保留参数检查点与训练记录；任务期状态更新不能代替部署期参数更新的证据 \\
$\mathcal V$ & 默认不在 Updatable 内（公理 1） & 系统同时能改实现与改考试，原分数不足以单独支持能力提高 & 纳入 Updatable 必须单列实验组并说明后果；还应检查候选能否改写评价所依赖的观测 \\
\bottomrule
\end{tabular}
\caption{Updatable每扩大一档要付什么。代价一列逐格对应表~\ref{tab:components}的更新代价性状；这里列出检验责任，不报告实测效果值。}
\label{tab:updatable}
\end{table}

竖着读这张表，应当看到成本来源的变化，而不是从低到高的固定排序。
一次很大的提示搜索可能比一次小范围技能修订更昂贵；经验记忆的索引方式与注入长度也会改变在线成本。
五类更新都需要与主张相匹配的检查；工件能否回退与外部效果能否恢复，还要分别记录。
公理 2 要求的不是把做过的检验再做一遍，而是每扩大一档就把六问的 Q6 重新答一次——
$P$ 的文件回退回答不了 $M$ 的噪声累积，$L$ 的入库验证回答不了 $G$ 的回归面。

% 先排完组件代价表，避免整页浮动表隔开方法例句。
\clearpage
式~\eqref{eq:ak} 的用处是给每个方法一个落点。
Reflexion~\cite{reflexion-2303} 落在 $M$：它把反思文本写进记忆，跨 episode 读取。
Voyager~\cite{voyager-2305} 的持久技能更新落在 $L$：它依据包含模型自评的成功判定存储程序，之后检索复用；这里的验证不等于独立测试断言。
DSPy~\cite{dspy-2310} 的编译可以优化模块中的示例等提示内容；原论文也讨论候选程序集成。应按具体编译配置登记 $P$，以及是否实际改变了 $G$，不能从框架名称推断每次都优化结构。
MEM1~\cite{mem1-2506} 的训练改变 $\theta$，部署时则由该策略生成压缩后的文本内部状态；两层变化需要分别记录。
这些方法的主要更新对象不同，成本、恢复条件与失效方式需要分别检查——这就是本书按组件而不是按论文组织材料的原因。

\section{把定义变成可复核的检查}

\subsection{重启之后，怎样确认经验真的留下来了}
\label{sec:ch01-reset-audit}

一个系统把上一轮的报错打印在屏幕上，另一个系统把报错写进文件，两者都可能在下一轮完全用不到这些信息。
保留信息和使用信息之间，还有读取路径、检索条件与版本选择。
检查一个自称会积累经验的系统时，可以先问一个很具体的问题：关掉本轮进程，再用新的任务目录启动，哪些文件会被主动加载？
这个问题把``模型好像记得''变成可检查的系统行为。

重启检查并不是要求每个方法都以磁盘文件保存经验。
数据库、远程索引或参数检查点都可以承担持久状态，只要能说明状态属于哪个版本、谁能读取、何时失效。
如果经验只存在于服务端持续交互状态，而客户端没有导出或重建机制，就应把服务端持续交互状态纳入系统边界。
此时能测试的是同一服务端状态下的持续行为，不能直接声称在任意重新部署后仍保留能力。

第一步是把初始版本完整固定下来。
所谓完整，不只是一份系统提示，还包括模型标识、工具接口、检索策略、持久库起点和验证规则。
只存提示而不存检索索引，下一轮加载的内容就可能已经改变；只存代码而不存模型服务的身份，同一程序也可能面对不同的求解器。
固定这些对象不是为了把系统做得复杂，而是为了让``版本改变''有明确的参照。

第二步是在若干来源任务后，保存候选更新与来源之间的连接。
每条经验至少能回答：由哪次任务产生、根据什么反馈提出、改动了哪个持久组件。
经验文本很短时也不能省去这些信息，否则后面发现一句规则错误，只知道删什么，不知道还需要复核哪些由它导出的条目。
这种连接属于可追溯性要求，并不自动证明经验是真的；来源本身仍可能出错。

第三步是清空任务内状态，再检查候选能否被加载。
工作上下文、临时文件和工具输出应按协议重置，持久物则按待检验版本显式载入。
如果这一步遗漏，所谓新任务可能仍在读取旧任务的中间答案。
一次测试全部通过，只能说明那个混合状态下任务能完成；无法分清收益来自持久更新还是未清空的工作目录。

第四步是把候选版本与冻结的旧版本分别送进相同的新任务条件。
两者要拥有同样的外部工具、预算和验收器，差异只放在本轮承诺检验的组件上。
这里先不讨论该用多少题或怎样估计区间，那是第\ref{ch03}章的工作；
本章要求的是，读者在看到任何分数以前，就能指出两条运行路径究竟在哪里不同。

这套检查也能得到有用的失败结论。
候选文件根本没有被加载，说明持久链路没有闭合；加载了却从未被检索到，说明存在读取条件问题；
被检索到并改变了行为，却在新任务上无益，才进入内容、适用范围与干扰的讨论。
把三种失败都写成``模型不擅长学习''，会让工程缺陷、选择缺陷和经验缺陷失去区别。

\subsection{文件哈希变了，能力就变了吗}

版本目录里多了一个日志时间戳，两个目录的哈希便不同，但求解器的行为可以完全相同。
因此 $A_{k+1}\ne A_k$ 必须相对于明确定义的组件与加载语义理解。
改动一个永远不会读取的注释文件，不等于改变任务策略；把同一条经验换个文件名，也不一定改变被检索的内容。
版本记录帮助证明工件不同，不能单独证明行为不同。

反过来，模型权重完全相同，也可能因为持久检索库变化而从不同起点作答。
一个系统增加了带来源的错误处理规则，另一个系统替换了常用工具的实现，这两种变化都不需要训练权重。
本书的五类组件是在系统层面划定的更新对象，不是把学习限定为梯度计算。
需要区分的是谁根据什么经验改变了它，而不是哪一种文件扩展名更像机器学习产物。

外部供应商更新了模型服务，会改变结果，却不能自动算作被测Agent自己的学习。
如果系统没有参与更新过程，也没有记录更新依据，那么实验应将其视作外部条件变化。
同一模型别名在不同日期返回不同后端，是版本解释的一种风险来源；仅保存别名不能排除它。
本书跨日期续跑的实验因此必须说明日期与身份限制，不能把所有差异都归给写入的经验。

人工编辑与自动更新也要在记录中分开。
工程师读完错误日志后改了一句提示，确实改变了系统版本，但它检验的是有人参与的更新流程。
若论文要研究Agent自主产生修改，就不能把人工筛选后的最佳提示当成自动修改器的原始输出。
人工帮助并不会使结果毫无价值；需要改变的是主张的对象、人工成本的记录，以及对应的比较基线。

\subsection{持久上下文、答案缓存与经验规则的边界}

一次持续交互的上下文可以跨越多次用户请求，因而``上下文''这个名字不能决定它属于哪一层。
如果协议把每次请求定义为新任务，而系统自动把上一请求的信息传给下一请求，相关状态就跨过了任务边界。
此时应把保留策略和读取过程写进系统版本的说明。
如果任务结束后上下文确实清空，它才是本章所说的临时工作状态。

答案缓存也会持久改变后续请求的执行路径。
缓存命中时，系统可以直接返回旧答案；若评价集重复出现原题，就可能显得越来越快、越来越准确。
这证明系统利用了重复请求，不自动证明它学到了可迁移的错误模式。
区分两者需要同时检查完全相同的请求、表述变化而需求相同的请求，以及要求发生变化的新请求。

以区间求和任务作教学例子，把某个输入对应的结果存下来，是一个输入到输出的记录。
把``确认题目是否包含右端点''保存为检查规则，则是在后续求解中提出一个条件判断。
后者看起来更一般，但它也可能在不相关任务中干扰求解；不能凭文本较抽象就直接授予迁移能力。
两类持久物都可以进入实验，只是比较任务必须与各自想证明的用途对应。

精确缓存尤其需要保留环境条件。
同一文件路径、同一句问题，在源文件、权限或接口版本变化后可能已经不对应同一个任务。
缓存键只取问题文本时，命中并不意味着旧答案仍有效。
这属于状态识别问题，与模型是否具有更强推理能力不同；提高模型规模也不会自动修复错误的缓存键。

经验规则还需要保留适用与不适用的区别。
``之前有一次边界错了''是一条观察，``以后所有循环都把上界加一''则把观察扩张成了规则。
后一个改动很容易降低来源任务的错误，却在右开区间本来正确的任务中制造新错。
本书关心的不是经验是否被压缩成一句话，而是压缩时保留了哪些决定条件，以及这些条件怎样在下一任务中被识别。

\subsection{沿着同一个错误，比较三种状态保留方式}

下面是一段教学推演，不包含新增实测结果。设有一个修复区间函数的任务，要求把闭区间内的整数相加。
初始实现误把右端点排除。Agent读出需求，修改循环边界，验收通过；随后收到另一个区间相关任务。
到这里为止，三种系统可以产生完全相同的第一条成功轨迹。区别要到任务结束后的处理才出现。

第一种系统只保留运行日志。日志中有需求、报错、修改与通过结果，但下一任务的启动配置没有日志读取器。
这种记录对研究者很有价值，能够帮助复现错误和核算成本；对求解策略而言，它仍是一份未参与决策的档案。
如果下一题也做对了，不能把成功归给这份档案。要支持这种归因，至少应先证明求解器确实获得了档案里的信息。

第二种系统在任务后保存一句经验：``区间题要检查左右端点是否包含''。
下一任务启动时，检索器按任务类型选中它，并把它加入提示。
此时持久更新与读取路径已经存在，但仍有两个不同问题：检索器是否把它送到了合适的位置，以及规则是否真的改善了解题过程。
一次检索命中回答不了后一个问题。它可能提醒了模型，也可能重复了模型本来就会做的检查，甚至挤掉更关键的需求说明。

第三种系统保存了完整的求和函数，下一任务又把这段函数作为参考代码读入。
如果新任务要求统计区间内满足条件的元素个数，旧函数里的端点处理也许有用，累加规则却不能照搬。
因此这个系统必须区分``检索到了代码''、``理解了适用条件''和``成功完成新任务''。
将完整代码放入上下文只建立了参考路径，并没有自动获得一个可直接调用且接口吻合的工具。

在三个系统中，验收第一题的动作完全相同，后续保留的对象却分别是无读取路径的日志、可检索的条件规则和供参考的程序。
它们需要不同的检验：日志要先查加载路径，规则要查条件保留和干扰，程序还要查接口、依赖与返回契约。
统一用``记住了经验''描述三者，会漏掉下一任务最可能发生的失败。

现在改变后续任务：题目明确采用右开区间。如果保存的是``把所有上界都加一''，来源任务的正确补丁便成了新任务的错误指令。
如果保存的是``先识别区间约定，再选择实现''，规则的适用范围更清楚，但仍要由任务中的语义识别来落实。
这里的差别不能仅靠经验文本的长度表达。长文本可能保留了条件，也可能堆入更多无关信息；短文本可能简洁，也可能删掉了决定正误的前提。

还可以改变接口而保留数学要求。例如下一任务接收一个包含左右边界的对象，而不是两个独立参数。
这会检验代码包装与接口适配，未必检验到更深的数学迁移。
若不把改变的轴写清楚，一次跨接口成功就可能被误说成跨领域泛化，一次接口失败也可能被误说成经验内容错误。
实验需要把新任务与来源任务的共同部分、变化部分分别列出。

最后改变预算条件。允许读取完整技能的系统，提示可能更长；从头求解的系统，则可能多走若干次工具调用。
如果只统计最终调用次数，前者似乎更省；如果计入技能建立、失败候选、检索和额外输入，比较可能发生变化。
本章不预先判断哪种方式更便宜，而要求在比较之前说明成本覆盖哪些阶段。
第\ref{ch08}章会进一步区分建立成本、复用成本与修订成本。

这个推演可直接转成一页实验说明，而无须先运行大量题目。
写明来源任务、持久对象、加载路径、新任务变化、对照版本与预算口径，就能发现设计中哪些问题还没有答案。
如果无法说明系统保留的是规则还是代码，应先补对象定义；如果无法说明新任务到底变了什么，应先补任务协议。
这比看到结果后再给成功起一个宽泛的名字更容易复核。

推演也说明为什么一次失败不必立即扩大更新权限。
接口不匹配可以先检查适配过程，规则错误可以先检查来源与准入，读取缺失可以先检查配置。
只有定位到现有更新对象确实不足，才有理由讨论是否开放新的组件。
否则，一个原本属于加载路径的小故障，可能被包装成需要同时训练模型、重构记忆和改写工具的大问题。


\subsection{评价边界：谁负责提出候选，谁负责决定采用}

同一个反馈信号可以有多种用途，但这些用途需要分别登记。
失败日志可以帮助提出修改，选择集结果可以帮助在多个候选中挑选，最终封存结果用来评价被选定的版本。
如果某个分数已经参与过候选选择，就不能再把它称作完全独立的最终评价。
问题不在文件是否被命名为``测试集''，而在信息是否已经流入修改过程。

验证器冻结并不等于Agent无法读到任何测试。
本书编程任务允许读取既有验收文件，重点限制其修改权限，并另外区分数据用途。
模型可以利用可见断言理解需求，这属于协议许可的反馈；模型把断言改成永真再获得通过，则改变了评价基准。
因此``可读''、``可写''与``可用于选择''是不同轴，不能用一个``封闭环境''概括掉。

只读文件保护也不是完备的安全证明。
它可以阻止工具直接重写指定文件，却不能自动排除依赖包变化、运行命令变化或其他环境通道对测试的影响。
本书当前装置的声明应限制在实现过的保护范围内；更强隔离需要额外的系统边界与检验。
第\ref{ch14}章讨论回归与恢复时，会把版本文件、运行状态和外部效果继续拆开。

提出候选和决定采用最好有不同的记录。
候选产生时，可以很大胆地尝试新的提示、记忆或工具；发布时则要说明通过了哪些条件，还有哪些风险未检查。
没有被发布的候选仍然是研究过程的一部分，它消耗的预算和暴露的失败模式都应留下。
只保存最终赢家，会让读者看到一个貌似便宜、稳定的更新过程，却看不到搜索曾经付出的代价。

\subsection{小改动为什么也需要发布清单}

持久记忆只增加一句话，也可能同时影响许多后续任务。
与单次答案不同，它的寿命超过当前交互，因此发现错误后的修复不能只改一次输出。
系统需要知道哪个版本开始读取这句话、哪些技能依赖它、哪些结果是在它存在时得到的。
本书把这些可检查的信息称为发布清单，它服务于复核与恢复，不是能力提升的评分表。

最小发布清单可以分成四组字段。
第一组识别更新前后的对象，包括组件、版本和内容校验信息；第二组指向来源轨迹与反馈；
第三组记录准入检查、预算和发布时刻；第四组说明恢复入口与尚未覆盖的条件。
字段不必依赖某种固定数据库，但必须能在结果被质疑时重新找到当时采用的工件。

\par\noindent\begin{minipage}{\linewidth}
下面给出一个教学用结构，它没有对应任何已执行的本书实验：
\begin{lstlisting}[language=Python]
release = {
    "component": "memory",
    "parent_version": "before-update",
    "candidate_version": "candidate-update",
    "origin_trace": "source-task-trace",
    "feedback_kind": "executable-test",
    "admission_record": "candidate-checks",
    "rollback_target": "before-update",
    "unverified_scope": ["new-task-family"],
}
\end{lstlisting}
\end{minipage}\par
这里用文字占位，避免把示意路径误当作真实工件。
实际运行必须把占位换成存在的对象，并能核对内容；无法找到来源的字段应记缺失，不能留一个看起来规范的路径充数。

有清单也不等于所有改动都能无损撤销。
提示文件通常能从旧版本恢复，模型参数如果保存了检查点也能回退；困难可能在存储、切换与相关状态的一致性。
但新版本已经发送的外部消息、已经写入其他系统的数据，不会因本地文件回退而自动消失。
因此本书把工件回退与效果补偿分开，不能用``参数更新不可逆''或``文本更新完全可逆''概括所有情形。

发布清单还支持一种朴素但重要的诊断：只移除最近的修改，观察问题是否还在。
如果同时改提示、检索、工具和模型，回退一项后结果没有恢复，也不能直接排除它的作用，因为组件之间可能存在交互。
这时应回到更小的对照设计，而不是不断给一个巨大的版本补标签。
版本粒度越大，追踪越省事，但归因与恢复通常需要处理的候选原因越多。

\subsection{从运行轨迹提出问题，不从单条轨迹宣布规律}

来源任务里的一个成功动作，可以帮助构造研究假设。
例如Agent先检查输入契约再修改代码，于是后来提出``保存契约检查顺序能减少相似错误''。
这条假设有明确的来源，但那个成功任务同时包含任务难度、模型采样和其他工具调用的影响。
从一条轨迹中选出自己喜欢的步骤，并不能证明去掉它就会失败。

失败轨迹也存在同样限制。
候选记忆被读入后任务失败，可能是记忆有害，也可能是任务本来更难，或验收环境已经发生变化。
如果只挑记忆最显眼的失败来叙述，就会让案例分析代替比较设计。
本书把案例用于定位机制、提出反例与解释执行链，把收益归因留给有明确比较对象的实验。

一个好的案例包需要保留区分竞争解释的入口。
除了任务陈述、候选改动和结果，还要告诉读者哪里可以看到未采用的版本，以及哪些解释仍未被排除。
例如``这条规则被加载过''可以由轨迹证明，``规则改善了后续任务''却需要另一种证据。
不把两种证据混为一谈，案例才不会变成给预设结论挑选的故事。

本书的主案例A刻意采用可执行的小任务，便于检查这个区别。
读者能看见初始缺陷、工具调用和验收判定，比只有自然语言评语更容易发现接口没有接上或结果被错误汇总。
代价是任务覆盖有限，失效模式经过人工选择；这种透明度不能代替真实项目的外部效度。
因此这里建立的是可复核的研究方法入口，不是对所有编程Agent能力的排行榜。

\subsection{阅读一份改进报告时，先找三个缺口}

第一类缺口是\textbf{对象缺口}：报告说更新了策略，却只展示一个任务最终输出的变化。
这时需要的补充不是更大的成功率表，而是持久组件在任务边界前后的清单和读取路径。
如果更新对象都没有被定位，后面的比较很难说明哪一种机制发生了作用。
对象缺口可以通过工程记录修补，不必先做新的大规模采样。

第二类缺口是\textbf{比较缺口}：候选获得更多模型调用、更长上下文或人工挑选，基线没有相应资源。
这类结果可能仍说明整个流程有效，却不能单独回答经验更新相对于增加推理预算的优势。
补充对照时应先指定想保持相同的资源和允许变化的资源，不能事后从多个成本指标里挑一个看起来对等的。
第\ref{ch03}章将把这种差别写成具体待估计量和预算账。

第三类缺口是\textbf{范围缺口}：来源任务通过，被写成新任务有效；有限格点有效，又被写成长期持续增长。
范围扩大通常需要新的检验条件，而不是更强的形容词。
同任务重试、族内变体、跨族任务、跨模型与跨环境，应各自说明状态继承与评价条件。
一个系统可以在其中一层有可靠证据，同时在另一层完全没有结果，这并不矛盾。

这三类缺口可以同时出现，但处理顺序不应颠倒。
先把对象和加载路径固定，再决定比较问题，最后扩大适用范围，能够减少在无法解释的装置上反复花费预算。
若反过来先跑大量题，再去猜系统究竟更新了哪里，更多样本只会把含混的过程记录得更精细。
本书从定义与实验台讲起，正是为了让后面的记忆、技能、工作流和参数更新共享这一套提问方式。


\section{六问：审视任何自进化系统的固定动作}
\label{sec:sixq}

\begin{enumerate}[label=Q\arabic*.]
  \item \textbf{改什么}：落在 $\theta,P,M,L,G$ 的哪一个，还是落在修改器本身。
  \item \textbf{根据什么改}：环境结果、执行测试、人类反馈，还是模型自评。主信号必须写明。
  \item \textbf{怎样改}：搜索、反思、程序生成、监督学习，还是强化学习。
  \item \textbf{什么时候改}：任务中、任务后、按周期，还是事件触发。
  \item \textbf{保留什么}：保留到 step、episode 还是 version 层。
  \item \textbf{怎样检验}：用选择验证集、回归集，还是封存集。
\end{enumerate}

每一问答不出来，都会以可预测的方式出错：
Q1 答不出，就可能把“上下文变长”误写成“持久组件已经更新”；
Q5 答不出，就会把一次任务内的重试说成跨任务进化；
Q6 答不出，升级和过拟合在数字上长得一模一样。

把六问用在 Voyager~\cite{voyager-2305} 上，得到的登记表是：

\begin{itemize}
  \item 改什么：$L$（技能库）。
  \item 根据什么改：环境反馈、解释器报错，以及另一 GPT-4 根据任务与环境状态给出的 self-verification 判断。
  \item 怎样改：模型生成程序，执行后把反馈回填；self-verification 判定成功后保存技能（原论文第2.3节）。
  \item 什么时候改：任务中遇到新目标、旧技能不够用时。
  \item 保留什么：version 层，技能跨任务复用。
  \item 怎样检验：去掉技能库的消融与新世界任务比较，为所报告设置中的组件作用提供证据；提示迭代配额不等于全部资源相等。
\end{itemize}

最后一行限制了前五行能支持的结论。Voyager 的第4节还明确承认自评可能误判，因此“技能已入库”不能直接改写成“独立验收已证明可复用”。
原论文的有无技能库对照值得保留；若进一步问同等总成本下是否优于增加重试，则还需另一种资源对齐设计。
本书把这些来源位置与主张边界登记在 \texttt{research/cards/}。

AFlow 提供了区分评价结果与成本口径的例子。AFlow~\cite{aflow-2410} 用蒙特卡洛树搜索改工作流 $G$ 与节点提示 $P$，
搜出的最优工作流跨任务保留，落在 version 层。它的检验有形式：验证集与测试集按 \cited{1:4} 划分，
每个候选工作流在验证集上跑 \cited{5} 次取均值，另外做了跨模型迁移表。
原论文图4明示成本是执行划分出的 HumanEval 测试集的总费用；这是测试执行口径，不能直接解释为包含离线搜索的总成本。
本书在该比较中没有定位到把搜索、验证与测试全部纳入的回本估计，也没有定位到总资源相等的扁平 best-of-N 对照。
这些缺口限制的是全流程收益主张，不否定其测试结果。表2已经比较跨模型使用的工作流：DeepSeek 所得工作流在 GPT-4o-mini 上低于为后者所得的工作流，却仍高于该模型的直接输入输出基线。
因此这里支持模型依赖性，不能写成作者未做迁移检验，更不能把相对模型专用版本的下降说成迁移后低于静态基线。

再看Reflexion、MEM1和JitRL，可以把运行期经验利用、策略训练和记忆操作的学习化分别定位。以下只讨论已核对的范围。

把六问用在 Reflexion~\cite{reflexion-2303} 上，可以把登记压缩为一句可检验的描述：不更新权重，把任务反馈转成反思，保存后供下一次尝试读取。
原论文第3节的反馈来源包含环境、启发式与模型评价，不能统一写成``没有外部接地''。
在本书登记中，还需补记一次尝试与一道新任务的关系；同题多次尝试与跨新题保留是不同主张。
判断适用范围时应检查具体任务协议，不能从一个环境中的反思失败推出整个框架没有持久更新。

MEM1~\cite{mem1-2506} 要分别登记训练与运行。
训练通过强化学习改变策略参数；运行时策略生成包含推理与摘要的文本内部状态，再把它带入后续交互。
原论文第3.1节区分了这两层对象，因此不能写成``每条任务记忆都已吸收到权重中''。
它说明了记忆管理策略可以被训练，却不单凭这一点证明部署后的每个新任务都会产生参数更新。
对本书而言，参数是否继续更新、文本状态何时清空、哪些状态跨越独立任务，都应分别列进版本协议。

把六问用在 JitRL~\cite{jitrl-2601} 上：

\begin{itemize}
  \item 改什么：运行时策略的 logit 偏置，不属于 $\theta,P,M,L,G$ 中的任何一类；同时 $M$ 作为非参经验记忆跨 episode 累积。
  \item 根据什么改：LLM评估器对完整轨迹生成逐步奖励，再折算为回报。
  \item 怎样改：对给定优势估计的KL惩罚目标求闭式解，再向 logit 加偏置，不更新基座参数。
  \item 什么时候改：每步检索经验并调制 logits，每个 episode 结束写入经验三元组。
  \item 保留什么：经验库 $M$ 在循环中累积，但要另查 episode 与任务的关系。原论文第5.2.1节对每个 WebArena 任务运行 \cited{5} 个 episode，不能把这五次尝试当作五道新题；算法1把记忆初始化放在 episode 循环之外，表5则另行检查只检索其他任务经验的冷启动设置。后者才直接涉及跨新任务的使用边界。
  \item 怎样检验：分别报告多次尝试的平均表现与末次表现，并做跨任务冷启动比较及“logit 更新对提示更新”的消融。WebArena-Lite 相对 WebRL 训练数据的留出，不等于 JitRL 在这些任务上没有再次尝试。
\end{itemize}

JitRL 值得单列，因为“任务策略如何更新”和“记忆操作是否被训练”是两个问题。
Q3 答“强化学习”只说明修改是怎么产生的，不说明修改了什么；JitRL 的强化学习学的是任务策略（选哪个动作），
记忆的写入、检索与价值估计全都是固定规则。于是六问在这里把两个常被混为一谈的问题分开了：
\textbf{有没有跨任务持久更新}需要 Q1 与 Q5 共同定位：JitRL 的经验库可以跨任务使用，这不由“无梯度更新”或 episode 名称单独决定；
\textbf{记忆操作有没有被学习}由 Q1 与 Q3 一起答，JitRL 的答案是“没有”（写入、检索、价值估计全是固定规则）。
第6章的记忆操作学习化讨论只问第二个问题，JitRL的记忆操作采用固定规则。

把这些方法的登记放在一起，六问的用途就清楚了：它把不同的变化对象与保留范围分别定位，帮助读者安排对应检验。
Reflexion 与 JitRL 都需要把记忆保留范围映射到具体任务边界；
Voyager 的技能库更新与 MEM1 的策略训练分别落在 $L$ 与 $\theta$，后者运行期的文本内部状态还需要另行记录。
差别不在方法名字上，在六问的答案上。

\section{主案例 A：一个反复犯同类错误的编程 Agent}
\label{sec:caseA}

全书需要一个能反复回访的场景，本书选用可执行测试的编程任务。

主案例A采用可执行验收的编程任务：Agent通过工具查看和修改文件，任务终判由执行验收给出。
各历史运行的工具循环并不相同，具体分派须按版本核对（附录B）；部分实验还向修改器提供模型反馈，这与任务终判是两条信号路径。
可执行测试提供了明确的验收入口，这是选它的第一个理由；测试是否符合题意、是否覆盖关键行为，仍然需要构造与核验。
第二个理由是任务文件与环境可以保存后重放，但模型输出、工具依赖与外部状态仍可能变化，复现性需要实际检查。

选这个案例还因为它同时满足三个技术条件。
验收入口明确：在给定测试下，通过与否不需要模型自评，但测试本身不能代替完整需求。
反馈可以帮助定位：失败测试和错误位置提供线索，却未必唯一确定根因或修复位置。
任务可复用：同一道题可以重复跑，配对比较与多次运行才有意义——
可重建的编程题便于固定输入并重复比较。若实际任务流持续变化，则应另外记录任务分布与出现顺序，
第13章讨论触发条件时会回到这一差别。

主案例 A 在本书里不是比喻。它的任务集已经建成：8 个任务族 $\times$ 8 题 $=64$ 题，
每题是一份带缺陷的小 Python 模块加一个独立验收测试
（协议见 \path{experiments/protocols/taskset-v1.md}）。

族不是按知识点分的，是按\textbf{失效模式}分的（表~\ref{tab:families}）。
分法取决于要问的问题：本书要问的不是“这道题会不会做”，
而是“在 A 类失效上积累的经验，迁到 B 类还剩多少”——本书按失效模式分族，使这个问题可操作化。按领域、仓库、工具或环境划分则对应其他迁移轴，应另定来源与目标范围。

\begin{table}[htbp]
\centering
\small
\begin{tabular}{@{}p{3.0cm}p{9.4cm}@{}}
\toprule
任务族 & 失效模式 \\
\midrule
\texttt{offbyone} & 边界与差一：区间端点含/不含、循环上界 \\
\texttt{mutabledefault} & 共享可变状态：可变默认参数与模块级容器的累积与别名 \\
\texttt{strindex} & 字符串索引与切片：方向、切分方式、大小写、端点字符 \\
\texttt{dictmiss} & 缺失键与计数初始化：默认值语义、累加被覆盖、缺失路径的返回 \\
\texttt{sortkey} & 排序键与稳定性：键取错字段、缺次级键、反向用错位置 \\
\texttt{errorcontract} & 输入契约：非法或越界输入应当返回什么，不是抛什么 \\
\texttt{classstate} & 对象状态累积：实例属性与类属性混淆、状态未更新、副作用与返回值不一致 \\
\texttt{recursionbase} & 递归基例与终止：基线取值、递归参数不收敛、嵌套未展开 \\
\bottomrule
\end{tabular}
\caption{主案例 A 的八个任务族。族内同构、族间可迁移检验是刻意的设计。}
\label{tab:families}
\end{table}

按失效模式分族要付一个代价：族内 8 题看起来像同一道题的八个变体。
这是有意的——族内变体可以检查较近的迁移，但不能替代跨族检验；
因此族内迁移的数字不能被单独解释成“泛化能力”。

64 题按族内编号分成四份，每份 16 题，每族恰好 2 题：训练、选择验证、回归、封存。
回归集与训练集由同一生成流程产生并使用不同题目，这控制了采集时间差异；它不排除模板重叠、答案泄漏或选择过程反复使用回归结果。
封存集不参与本轮候选产生与选择；终评查询要留痕，不能把查询后调整的版本继续称为同一次独立终评。

一道题比一段描述更能说明这类任务在测什么。以 \texttt{offbyone-01} 为例，它的题面是：
实现 \texttt{sum\_inclusive(a, b)}，返回从 \texttt{a} 到 \texttt{b} 之间所有整数的和。初始代码只有两行：

\begin{lstlisting}[language=Python]
def sum_inclusive(a, b):
    return sum(range(a, b))
\end{lstlisting}

缺陷出在 \texttt{range}：它是右开区间，\texttt{b} 从来没被加进去。
验收测试不检查实现，只断言三个字面量期望值（\texttt{test\_main.py}）——
\texttt{sum\_inclusive(1, 5)} 应为 \calcnum{15}，缺陷代码给出 \calcnum{10}；
\texttt{sum\_inclusive(3, 3)} 应为 \calcnum{3}，缺陷代码给出 \calcnum{0}；
\texttt{sum\_inclusive(-2, 2)} 应为 \calcnum{0}，缺陷代码给出 \calcnum{-2}。
三条断言在初始代码上全部失败；这一点读代码就能算出来，不必运行任何 Agent。
构建器另外把含缺陷版本至少失败一条断言写成了强制条件，作为这道题能被使用的前提。

它归在 \texttt{offbyone} 族，是因为三个用例测的是同一件事：区间右端点含还是不含。
第二个用例把区间压到两端相等的退化情形，第三个让它跨过零点——
两个用例都在把同一个缺陷推到藏不住的地方，这正是族的定义所要求的。

同一道题也演示了 $\mathcal V$ 为什么必须被挡在 Updatable 之外。验收测试里没有参考实现，只有期望值：
沙箱内的文件对 Agent 可读，测试里一旦带一份正确实现，任务就退化成复制。
\texttt{offbyone-01} 的 \texttt{test.json} 把 \texttt{test\_main.py} 列进 \texttt{protected}，
通过指定写工具直接改该测试的请求会被拒绝。这个工具级规则有助于区分实现修改和测试修改，但没有把宿主文件真正设成所有候选进程都不可写。查看历史成绩时，应把现有路径保护与尚未实现的强进程隔离分开，不用\texttt{protected}这个字段名代替权限验证。

族内同构的用途在这道题上也能直接看出：在 \texttt{offbyone-01} 上学会检查右端点的系统，
遇到同族的另一题是否还能用，正是族内差异被刻意压小后要检验的问题。
反过来，在 \texttt{offbyone} 上积累的经验迁到 \texttt{errorcontract}
（非法输入该返回什么，不是抛什么）还剩多少，是完全不同的问题。
第 5 章的跨族迁移问的是后一个。前一个还要分清：在同一道题上重复运行主要帮助估计运行波动，
在同族新题上测试则可考察族内迁移；两者都不能自动替代对新任务族的检验。

本章用这套任务集定义可检查的对象。能力读数从第3章起展开，并绑定实验与运行标识。

% ---------- 4. 对照实验 ----------
\section{预算比较需要怎样的对手}
\label{sec:firstcontrast}

本章开篇那组数字的完整形态见表~\ref{tab:harness}。请注意它的比较方式：
不仅问“演化有没有让分数变高”，还问\textbf{在已声明的采样机会约束下，不同分配方式表现如何}。
机会约束与完整资源账是两层问题，不能只凭表中分数回答后一层。无单元测试时，并行臂由模型自评从候选轨迹中选出最终提交；这里的pass@1评价选出的提交，不是基础模型只生成一次的成绩。原文结果取\cited{2}次独立运行的均值，点估计排序仍须与差异的不确定性分开阅读。

\begin{table}[htbp]
\centering
\small
\begin{tabular}{@{}p{3.4cm}cccc@{}}
\toprule
方法 & \multicolumn{3}{c}{pass@1} & 平均 \\
\cmidrule(lr){2-4}
 & Opus 4.6 & GPT-5.4 & GPT-5.4 mini & \\
\midrule
初始 harness 直接采样 & 69.9 & 75.3 & 59.4 & \textbf{68.2} \\
并行采样 & 74.7 & 79.2 & 62.9 & \textbf{72.3} \\
顺序精化 & 73.0 & 73.0 & 61.8 & \textbf{69.3} \\
harness evolution & 71.4 & 69.7 & 61.3 & \textbf{67.4} \\
harness scaling & 76.0 & 78.1 & 61.2 & \textbf{71.8} \\
\bottomrule
\end{tabular}
\caption{Terminal-Bench 2.1，无单元测试反馈，pass@1。数字均为原论文报告（\cite{harness-eval-2607}，Table 1），非本书实测。最后一列加粗便于比较，非原文格式。}
\label{tab:harness}
\end{table}

所用v2的 Table 2 另报有单元测试条件下的结果：pass@1 平均，并行采样 \cited{86.0}，harness evolution \cited{75.8}；
pass@5 平均，顺序精化 \cited{91.8}，harness evolution \cited{86.2}。该表只含两个模型，
且测试同时用于反馈和候选选择；不能把它与表~\ref{tab:harness} 的三模型均值相减，解释成单纯增强反馈的因果效应。

\paragraph{版本改变了哪些证据。}
这里保留v2数值，使历史案例的模型与选择规则可以复核。2026年10月5日的v4在Table~2和Table~3增加GPT-5.4 mini；Table~3的三模型平均增量为\cited{0.0}个百分点，不能与v2的两模型均值\cited{0.6}混用。v4 §5还报告长时游戏中的任务内harness scaling收益：以环境持续反馈改进当前游戏，并记录整段运行的最好成绩\cite{harness-eval-2607-v4}。这补充了“什么条件下值得更新”的正向证据；它检验的是当前任务适配，不能替代新任务的留出检验。

\par\noindent\begin{minipage}{\linewidth}
\begin{proposition}[P1]
在什么任务与预算条件下，把资源用于跨任务持久更新，比用于任务内重试更有效？
\end{proposition}
\falsifier{须先使比较能够识别同预算下的分配效应。有效证据若支持相反方向，可反对对应范围的优越性主张；若区间上界低于事先要求的最小收益，可排除该幅度的收益。原 3 个百分点设计目标保留，但区间跨零、点估计小于目标或精度不足都不自动证伪。原有方向性假设另存于附录\ref{sec:appg-history}。}
\end{minipage}\par

\par\noindent\begin{minipage}{\linewidth}
\begin{proposition}[P1b]
共享同一起始记忆后，继续更新比只读复用多带来什么？
\end{proposition}
\falsifier{归因于更新本身时，两臂须共享起始记忆并对齐信息、选择和成本范围。有效比较支持相反方向才构成相应反证；区间包含 0 表示结果尚未排除零差异，也可能容许正负效应，不能据此宣布更新无益。这些条件界定所问的增量，而不预设更新必然有益。}
\end{minipage}\par

表~\ref{tab:harness} 挑战了“留下修改就会胜过增加采样”的简单推断，但它并不直接完成本书 P1 的跨任务比较。
后面各章讨论收益时，需要设置适合该任务的重试与候选选择对照，否则增益可能来自不同的预算分配。
它也提示 P1b 所需的对照：把同一份持久物保持只读，再与允许更新的版本比较，
才能把起始信息与后续更新分开。这是本书的设计要求，不是对全部文献作“此前无人检验”的判断。
历史 E01 已有只读格，但其起始经验来源与更新格不同，尚未完成这项识别（第 3 章）。

E01 的四格按持久范围组织。历史运行对齐 episode 配额，但并未把完整成本与所有执行机制对齐。

\begin{table}[htbp]
\centering
\small
\begin{tabular}{@{}p{1.2cm}p{3.6cm}p{4.6cm}p{3.4cm}@{}}
\toprule
格 & 预算花在 & 系统有没有持久物 & 关联问题 \\
\midrule
格 1 & 静态执行，重复取均值 & 无 & P1 的对手 \\
格 2 & 建了持久物，只读不写 & 有，但不生长 & P1b 的对手 \\
格 3 & 每个任务单独适配，用完即弃 & 有，但不跨任务 & 与格 4 对照 \\
格 4 & 跨任务持久更新 & 有，并供后续任务使用 & P1 \\
\bottomrule
\end{tabular}
\caption{E01 的四格及对照目的。历史运行共享 episode 配额，不等于共享相同总调用或总费用；执行循环与经验来源的差异见第 3 章审计。}
\label{tab:grid}
\end{table}

格 2 的用途是让“拥有可读取的信息”和“在线改变这些信息”分别进入比较。
它可能优于无记忆系统，也可能因检索失败而没有差别；不能预先把其成绩解释为持续更新的收益。
只给一个臂起名为“只读基线”仍不够，还要检查它读的是否确实是另一臂的同源初始材料。

任务集定义了训练、选择验证、回归与封存四个各 16 题的分区，这是可用任务的结构，
不是四个分区都已被 E01 完整执行的证明。历史 E01 每格训练 episode 配额为 8，
封存阶段使用 16 题；具体重复、调用上限和流程见第 3 章。目录存在不能替代实际查询日志。

更新格额外发生的写入、检索和重写，必须连同准备材料的成本一起入账。
若只对齐做题次数而忽略这些费用，比较只能称为相同 episode 配额下的方案比较。
E01 新版据此准备统一资源账与共同起始记忆合同；其离线检查只说明部分控制逻辑已实现，
尚未形成新的等预算能力结果，更没有使历史结果获得新的因果含义。

\section{换任务、换模型、换反馈器之后}
\label{sec:conditions}

一个结论要说明适用条件。下面三种变化提示不同的外推风险，它们并不都由同一个纯干预实验识别。

\textbf{换任务}。前面已经看到：held-out 上平均只涨 \cited{0.6} 个百分点，GPT-5.4 上为 \cited{0.0}。
原作者把原因归结为修改编码了任务特定的捷径。这句话和第 5 章的\textbf{拟合目标}是同一件事的两种说法：
修改可能利用了当前任务与求解流程的特征。不过，换任务后的成绩变化不能单独证明问题来自读取者；
要区分任务分布、材料内容和模型读取方式，还需要相应的交叉对照。

\textbf{换模型}。表~\ref{tab:harness} 里 harness evolution 在 GPT-5.4 上相对直接采样是倒退的（\cited{69.7} 对 \cited{75.3}），
在 Claude Opus 4.6 上是前进的（\cited{71.4} 对 \cited{69.9}）。
这是同一方法在不同模型配置下的结果，不等于把完全相同的更新产物交给两个模型读取。
因此它提示模型依赖，却不能单独估计产物跨模型迁移的损失。AFlow~\cite{aflow-2410} 的跨模型工作流比较
提供了另一种检查思路，具体搜索与执行模型、基线和成本范围见第 9 章。
第 5 章据此区分“分别为每个模型优化”和“固定产物更换读取者”，两者回答的问题不同。

\textbf{换反馈器}。单元测试可以同时改变求解期间的修正和终点的候选选择。
前述外部比较把这两条路径同时打开，不能仅凭总分变化把收益全部归给反馈质量。
第 4 章的 P2 需要在共同执行流程里交叉操纵反馈与更新对象；跨表均值或跨论文排序都不能替代这一设计。

换反馈器还会改变验证依据及其覆盖范围。表~\ref{tab:verif}按信号来源列出四类；类别有助于提出检查问题，不能单凭层号给具体系统的可靠性排序。

\begin{table}[htbp]
\centering
\small
\begin{tabular}{@{}p{1.0cm}p{3.0cm}p{5.0cm}p{4.2cm}@{}}
\toprule
层 & 信号 & 性质 & 例子 \\
\midrule
1 & 形式验证器 & 结论依赖规格、证明规则与检查器 & 证明检查器、类型系统 \\
2 & 执行反馈 & 覆盖有限；反复针对固定检查优化可能失真 & 测试、编译器、基准 \\
3 & 学习式评审 & 取决于证据、工具与核查流程，也可能成为优化目标 & reward model、LLM-as-judge \\
4 & 内在信号 & 与外部目标的对应关系需验证；成本随实现变化 & 置信度、自一致性、似然 \\
\bottomrule
\end{tabular}
\caption{反馈的可靠度四级（\cite{rsi-survey-2607}，\S 5.2）。原论文自陈“自改进强度沿此层级递减”是语料中的
定性观察，\textbf{不是测量定律}；第 4 章报告 E02 已测操作的效应与限制，其 formal 臂实际为单元测试，不能据此声称完整检验了形式验证层级。}
\label{tab:verif}
\end{table}

对第 1 章有用的只有一句：换反馈器不只是“变强一点”或“变弱一点”。
把候选生成阶段的执行反馈改成内在信号，会改变优化依据，却不必同时改变最终验收。
只有把该内在信号也用作最终效果的唯一裁判，才把原来不同的职责合并了。公理1要求明确并冻结评价边界，不能把所有内部反馈变化都说成违反这条约定。

RuHarness的v1给出了一个具体例子\cite{ruharness-2610}：它同时维护运行脚手架和诊断用的行为判据，并随执行轨迹更新判据及其优先级。论文规定，判据打标时不读取最终任务结果；其后的判据准入与效用统计仍使用这些结果。候选脚手架按任务分数在单独的选择验证集上比较，最终报告另用封存集。这一区分保留了可变诊断依据、候选选择和最终评价各自的用途；它不表示整个更新过程都不看任务结果，也不能单凭诊断分上升判断任务能力提高。

% ---------- 6. 选择原则 ----------
\section{把机制选择变成可检验的设计}
\label{ch:why}

方法谱系帮助读者找到候选路线；采用其中一种，还需要把机制接到具体任务、资源与评价上。本书的主线是完成这种连接：先指出实际变化，再建立对应比较，最后沿工件解释结果。记忆、提示、技能、工作流和参数学习由此共享一种提问方式，同时保留各自的实现与成本特点。

已有综述为这条路径提供了不同入口。表~\ref{tab:diff}列出本书采用的框架，以及读者如何把它们写进自己的设计。

\begin{table}[htbp]
\centering\small
\begin{tabular}{@{}p{3.0cm}p{5.0cm}p{5.2cm}@{}}
\toprule
来源 & 使用的框架 & 接到实验中的对象 \\
\midrule
\cite{self-evolving-survey-2507} & 四问分类与保持、迁移、效率指标（§7.1）；比较条件（§7.3.2） & 目标任务分布、预算范围和逐次记录 \\
\cite{self-evolving-comprehensive-2508} & 四组件反馈环；搜索空间与优化算法$(\mathcal S,H)$（§3.1、§3.5） & 允许修改的对象、候选产生与停止规则 \\
\cite{self-improvements-survey-2607} & $A_t=(\theta_t,\Sigma_t)$、累计预算与报告要求（§8） & 轨迹、资源、保持结果和评价者身份 \\
\cite{rsi-survey-2607} & 两轴分类、验证层级、有界与开放边界（§5、§7） & 操作定义、对应对照和适用范围 \\
\cite{coding-survey-2608} & v3的持久目标、时机与信号；评价指标（§5.2） & 准入、退役、恢复及复用成本的观测量 \\
\bottomrule
\end{tabular}
\caption{从综述框架到实验对象。右列说明本书的使用方式，具体结果仍沿相应原始来源核对。}
\label{tab:diff}
\end{table}

社区清单也能帮助发现材料，例如\texttt{lobehub/awesome-rsi}\cite{awesome-rsi}。使用这些入口时，先记录方法与问题，再返回原始论文、代码或运行记录确认具体机制。这样，阅读路线和结论依据就能分别定位。

\subsection{先明确资源如何分配}

设一个团队已经有了可执行的静态Agent，准备投入更多资源。它可以增加同题尝试，也可以把来源经验写入持久物。要比较这两种选择，应同时固定任务范围、可见反馈、起始材料和候选选择规则，并说明共同约束的是调用、token还是墙钟。E01四格将这些选择摆在同一个表中；第3章沿历史实现检查各格实际获得了什么。

只读持久对照在这里尤其有用。让两臂从同一份记忆开始，一臂保持只读，一臂允许继续写入，就能把初始信息与持续更新分别放进设计。随后还需核对写入和选择所用资源，比较才回答预定问题。这个过程把“加了记忆有用吗”推进到一个更具体的问题：在给定信息和资源下，持续更新增加了什么？

\subsection{再明确成本怎样累积}

技能可能把发现阶段的工作搬到部署之前，记忆则可能让每个后续任务多读一段上下文。因此成本必须沿时间展开：准备、构建、使用、维护和验证分别由什么事件产生，共享准备如何计一次，失败尝试保存在哪一笔账中。第8章的成本模型和附录F的历史账目据此连接起来。

一个最小教学比较可以写成两条直线：从头求解每次花$c_0$，技能先花$C_{\mathrm{build}}$，以后每次花$c_1$。在$c_0>c_1$且质量满足共同要求时，连续交点为$C_{\mathrm{build}}/(c_0-c_1)$。加入维护或改变任务分布后，应按实际分段重新计算。读者由此能够指出回本判断需要哪几个输入，而不是从一次调用减少推断长期收益。

\subsection{让观察结果能够改变设计}

每项比较都应事先说明哪些观察支持采用，哪些会要求修改方案，哪些仍不足以区分解释。平均收益、保持、迁移和成本可能给出不同方向；工程决策需要清楚的质量约束，科学判断需要与主张匹配的干预和不确定性。第16章把这些条件组织为十个研究问题，并列出现有观察与下一步设计。

本版的材料路径见表~\ref{tab:three}。读者可以从一个问题进入相应章节，再沿归档复核它，而无须先接受一条预定的成功故事。

\begin{table}[htbp]
\centering\small
\begin{tabular}{@{}p{2.6cm}p{4.8cm}p{5.8cm}@{}}
\toprule
问题 & 工件与阅读入口 & 可以跟随的动作 \\
\midrule
资源怎样分配 & 第3章E01四格与对应轨迹 & 对齐尝试、初始材料、选择过程与费用范围 \\
成本怎样累积 & 第8章模型、第15章E05v2案例、附录F & 从实际事件合账，同时核对任务完成质量 \\
主张怎样判断 & 第5章E03审计、第16章研究问题表 & 先确认有效记录集，再计算对应差值并限定解释 \\
\bottomrule
\end{tabular}
\caption{从问题到工件的三条阅读路线。各运行的已完成范围和待检验条件在对应章节列明。}
\label{tab:three}
\end{table}

% Shared evidence-chain figure for chapter 1.
\begin{figure}[htbp]
\centering
% Evidence responsibilities; no literature ranking or claim of an empty field.
\begin{tikzpicture}[
  box/.style={draw=gray!65,rounded corners=2pt,align=left,text width=3.55cm,
              minimum height=2.8cm,inner sep=8pt,font=\small},
  arrow/.style={-{Stealth[length=2mm]},thick,gray!70}
]
\node[box,fill=blue!4] (source) at (0,0) {\textbf{来源与设计}\\[4pt]
术语与机制定位\\比较对象与目标量\\预算、信息和判据};
\node[box,fill=green!4] (record) at (4.8,0) {\textbf{实现与观察}\\[4pt]
版本和执行权限\\逐次记录与失败账\\离线检查和实测分开};
\node[box,fill=orange!5] (claim) at (9.6,0) {\textbf{证据状态与范围}\\[4pt]
可识别性与不确定性\\结果支持到哪一步\\未解决项继续保留};
\draw[arrow] (source.east) -- (record.west);
\draw[arrow] (record.east) -- (claim.west);
\end{tikzpicture}

\caption{从来源到结论的三层责任。图示为本书组织方法，不是对已有文献的穷尽排名。完成设计或软件检查，不会自动产生有效能力比较；第16章据此组织研究问题与后续比较。}
\label{fig:survey-axes}
\end{figure}


图~\ref{fig:survey-axes}把来源、执行与判断连成证据链。六问给出对象，运行记录连接实际动作，比较设计决定差值回答什么；这条链也是读者组织下一份实验说明的起点。

\section{建立证据边界与阅读路线}
\label{sec:boundaries}

全书将三类材料分别定位。外部论文和技术报告保留其任务、模型、指标、预算与选择条件，重要主张返回原始来源核对。来源的机构名称不能代替方法审查；公开元数据确认身份，具体段落和结果才支持具体命题。

本书运行沿配置、逐次轨迹、汇总与解释读取。来源与汇总一致时，还要检查对照能否识别所问作用；记录冲突时，则先确定可追溯的分析记录集。第5章对E03基线的处理展示了这一顺序，研究问题及已有观察集中整理在第16章，历史假设另存于附录\ref{sec:appg-history}。

教学计算写明输入和假设，用于展开一个关系如何成立、改变条件后怎样变化。读者可以用同一公式代入自己的测量，也应保留新输入的身份。由此，机制解释、模型推算和实际观察能够连成论证，每一项又都有明确的复核入口。

\begin{principlebox}
判断一个系统是否真的在自进化，只看一件事：\textbf{后续任务的出发点是否被先前的经验改变，且这个改变可以被独立检验}。
出发点没变，不能据此声称跨任务持久更新；改变已经发生而效果证据不足，应记录为效果未明。是否属于重试、成本是否更高，仍要看实际过程。
\end{principlebox}

\begin{pitfallbox}
\textbf{反例：把发生更新误当成收益保证。}
表~\ref{tab:harness} 里，harness evolution 在无单测条件下低于不做演化的起点（\cited{67.4} 对 \cited{68.2}），
而并行采样的三模型平均点估计最高（\cited{72.3}）。
\textbf{反例：把过拟合包装成迁移。}
同一批 harness 改进迁到 held-out 任务上，平均只涨 \cited{0.6} 个百分点，GPT-5.4 上为 \cited{0.0}。
\textbf{反例：把提示变长当作持久更新的证据。}
本章判据要求先检查保留边界：若上下文只在任务内增长、结束后清空，它不改变下一任务的起点。
\end{pitfallbox}

\section*{习题}
\begin{enumerate}[label=1-\arabic*.,ref=1-\arabic*]
  \item\label{ex:ch01-1} 找一个你熟悉的 Agent 系统，按式~\eqref{eq:ak} 指出它哪些组件在 version 层发生变化，哪些只在 episode 层。
  \item\label{ex:ch01-2} 表~\ref{tab:harness} 中五个方法的均值能支持哪些描述，尚不能支持哪些推断？选定一个主要比较，说明还需哪些逐题结果、重复运行与候选选择记录，以及怎样估计差异的不确定性。
  \item\label{ex:ch01-3} 若一个系统的 $\mathcal V$ 本身可被修改，本书的哪条公理会被触及？在什么条件下这种修改仍然可辩护？
  \item\label{ex:ch01-4} 举一个“看起来在进化、实际只是重试”的例子，并说明你如何设计对照把两者分开。
\end{enumerate}

\section*{常见误解}
\begin{itemize}
  \item \textbf{“有记忆就是自进化”。}有持久物不等于根据新经验继续更新。E01 的“只读不写”格专门用来分开这两件事。
  \item \textbf{“分数涨了就是学会了”。}可能来自更多重试、更简单的任务分布，或者选择验证集被反复查询。
  \item \textbf{“用了强化学习就是自进化”。}要先确定更新对象与保留边界。第6章的 JitRL 还说明，任务策略由强化学习调制，不等于记忆操作本身也经过训练；操作规则固定也不否认经验库发生持久变化。
\end{itemize}

\clearpage
\section*{进一步阅读}
\begin{itemize}[itemsep=2pt,topsep=4pt,parsep=0pt]
  \item \cite{harness-eval-2607}：比较rollout约束下的任务内改进与独立留出题表现，阅读时区分两套实验及其对照。
  \item \cite{rsi-survey-2607}：两轴分类、反馈条件与长期评价问题，为本书从机制走向测量提供参照。
  \item \cite{react-2210}：自进化之前的执行基线，本书的对照层起点。
  \item \cite{self-correct-negative-2310}：缺少外部反馈时的推理自我修正实验，第4章结合其具体设置讨论反馈质量。
  \item \cite{coding-survey-2608}：六类持久目标与六类信号的来源，第 4、8 章的取材处。
\end{itemize}
